%% ============================================================
%%  A Click is All You Need — 汇报讲稿 Speaker Script
%%  总时长：约 10 分钟 | Engine: XeLaTeX | Compile twice
%% ============================================================
\documentclass[11pt,a4paper]{article}

\usepackage[UTF8,scheme=plain]{ctex}
\setCJKmainfont{Kaiti SC}
\setmainfont{Times New Roman}

\usepackage[top=2cm,bottom=2cm,left=2.5cm,right=2.5cm]{geometry}
\usepackage{xcolor,parskip,fancyhdr,enumitem}

\definecolor{speakblue}{RGB}{0,52,102}
\definecolor{notegray}{RGB}{80,80,80}
\definecolor{speakbg}{RGB}{232,242,255}
\definecolor{notebg}{RGB}{245,245,245}
\definecolor{gold}{RGB}{180,140,40}
\definecolor{accent}{RGB}{196,72,49}

\pagestyle{fancy}
\fancyhf{}
\fancyhead[L]{\small\color{speakblue}\textbf{A Click is All You Need} — 汇报讲稿}
\fancyhead[R]{\small\color{speakblue}KyrieWong7 · April 2026}
\fancyfoot[C]{\small\thepage}
\renewcommand{\headrulewidth}{0.4pt}

%% 幻灯片节标题
\newcommand{\slidesec}[2]{%
  \vspace{14pt}
  \noindent\rule{\textwidth}{0.8pt}\\[-2pt]
  {\color{speakblue}\large\bfseries 幻灯片 #1 \quad #2}\\[-4pt]
  \rule{\textwidth}{0.4pt}
  \vspace{4pt}}

%% 口播内容框
\newcommand{\speakbox}[1]{%
  \vspace{4pt}
  \noindent\fcolorbox{speakblue}{speakbg}{%
    \begin{minipage}{\dimexpr\textwidth-2\fboxsep-2\fboxrule}
      \vspace{3pt}
      {\small\color{speakblue}\bfseries 🎤 口播内容}\\[3pt]
      {\normalsize #1}
      \vspace{3pt}
    \end{minipage}}
  \vspace{4pt}}

%% 舞台提示框
\newcommand{\notebox}[1]{%
  \vspace{3pt}
  \noindent\fcolorbox{notegray}{notebg}{%
    \begin{minipage}{\dimexpr\textwidth-2\fboxsep-2\fboxrule}
      \vspace{2pt}
      {\small\color{notegray}\bfseries 📋 提示 / Stage Direction}\\[2pt]
      {\small\color{notegray} #1}
      \vspace{2pt}
    \end{minipage}}
  \vspace{3pt}}

%% ============================================================
\begin{document}

\begin{center}
  {\LARGE\bfseries\color{speakblue} A Click is All You Need}\\[6pt]
  {\large 研究生PPT汇报解救器 · 汇报讲稿}\\[4pt]
  {\normalsize KyrieWong7 \quad|\quad April 2026 \quad|\quad 约 10 分钟}\\[4pt]
  \textcolor{gold}{\rule{0.5\textwidth}{0.6pt}}\\[4pt]
  {\small\color{notegray}
    本讲稿由本项目工作流自动生成，与 \texttt{slides.pdf} 逐页对应。\\
    \textit{This script was generated by the very pipeline it describes.}}
\end{center}

\vspace{10pt}

%%
%% SLIDE 1 — COVER
%%
\slidesec{1}{封面 Cover}

\notebox{【展示封面幻灯片。调整麦克风，等教室安静下来。】}

\speakbox{%
大家好。今天我来分享的是一个叫做 \textbf{「A Click is All You Need」} 的项目——一个让你在 Cursor IDE 里对 AI 说一句话，就能拿到专业学术汇报 PDF 的工作流。

\vspace{4pt}
Good morning / afternoon, everyone. Today I'll be presenting \textbf{"A Click is All You Need"} — a Cursor-native workflow that turns one sentence into a professional, bilingual academic presentation.

\vspace{4pt}
顺便说一句：这套幻灯片本身，就是用它自己的流水线生成的。
\textit{By the way: these slides were generated by the very pipeline I'm about to describe.}
}

\notebox{【停顿 1 秒，让最后那句话沉淀一下。预计用时：30 秒】}

%%
%% SLIDE 2 — HOOK
%%
\slidesec{2}{钩子：深夜的 Deadline}

\notebox{【翻到 Slide 2。指向左侧 Chiikawa 图。】}

\speakbox{%
我想先从一个大家都很熟悉的场景开始。

\vspace{4pt}
11 点了，明天 9 点有个学术汇报，PPT 还没做完。你打开 PowerPoint，发现字体在 Mac 上好好的，昨天在实验室的 Windows 上全乱掉了。你找了半小时适合的配图，但分辨率不够，嵌进去一片模糊。更糟的是——PPT 算是做完了，但讲稿，你根本没开始写。

\vspace{4pt}
This is the universal graduate student nightmare. And it happens every single time.
}

\speakbox{%
这不是个人能力的问题。这是工具的问题。

\vspace{4pt}
PowerPoint 是为商业演示设计的，不是为学术汇报设计的。它让你花大量时间在排版上，而不是在思想上。

\vspace{4pt}
\textbf{「A Click is All You Need」} 就是为了解决这个问题而生的。
}

\notebox{【预计用时：1 分 30 秒】}

%%
%% SLIDE 3 — ROADMAP
%%
\slidesec{3}{全局路线图 Roadmap}

\notebox{【翻到 Slide 3。指向 TikZ 四节点流程图。】}

\speakbox{%
今天的汇报分四个部分。

第一节，问题诊断——我们为什么需要这个工具？
第二节，解法流程——8 阶段全自动流水线是怎么工作的？
第三节，技术规范——字体、配色、用户配置怎么管理？
第四节，效果对比——Before 和 After 的差距有多大？

\vspace{4pt}
Four sections: The Problem → The Pipeline → Tech Spec → Results.
}

\notebox{【预计用时：30 秒】}

%%
%% SLIDE 4 — THE PROBLEM
%%
\slidesec{4}{S1 三大 PPT 痛点}

\notebox{【翻到 Slide 4。】}

\speakbox{%
我把传统学术 PPT 制作的痛点归纳为三类。

\vspace{4pt}
\textbf{第一，排版失控。}字体跨平台乱版，公式截图嵌入模糊，手动逐页调字号，效率极低。

\vspace{4pt}
\textbf{第二，内容分散。}PPT、讲稿、参考文献在三个不同文件里，改一处要同步三个地方，版本混乱。

\vspace{4pt}
\textbf{第三，质量无保证。}做完就提交，问题在讲台上才发现——文字被裁掉，图表节点重叠，叙事逻辑跳跃，听众跟不上。

\vspace{4pt}
研究生平均每次汇报花 5 到 8 小时，其中大量时间都在和工具做无谓的搏斗。
}

\notebox{【预计用时：1 分钟】}

%%
%% SLIDE 5 — 8-STAGE PIPELINE
%%
\slidesec{5}{S2 解法：8 阶段全自动流水线}

\notebox{【翻到 Slide 5。指向左侧 TikZ 竖向流程图。】}

\speakbox{%
我们的解法是一条 8 阶段的全自动流水线，完全运行在 Cursor IDE 内部。

\vspace{4pt}
整个流程是这样的：首先，AI 摄入你的素材——PDF、DOCX 或者你的口头描述；然后搭建脚手架，从模板复制项目结构；接下来逐页生成双语内容；同步规划并生成配图；然后 XeLaTeX 编译两遍；再由三个专责 AI agent 并行审查；审查通过后自动生成讲稿；最后交付验证。

\vspace{4pt}
输入：一句话 + 你的素材。输出：slides.pdf 和 script.pdf，两个文件，缺一不可。
}

\notebox{【预计用时：1 分钟】}

%%
%% SLIDE 6 — WHY NOT OFFICE?
%%
\slidesec{6}{为什么不用 Office？}

\notebox{【翻到 Slide 6。指向对比表格。】}

\speakbox{%
这是一张我很喜欢的对比表。八个维度，左边是 A Click，右边是 Office 或普通模板。

\vspace{4pt}
排版精度：XeLaTeX 是 pt 级别的精确，Office 是你用鼠标拖到看起来差不多。

字体一致性：A Click 改一行 config.tex 全局生效，Office 逐页手动点。

跨平台：PDF 是 PDF，不存在"你那里好看我这里乱"的问题。

图形质量：TikZ 是矢量图，任意放大不失真；截图嵌入在投影仪上会糊。

版本控制：.tex 文件可以 git diff，.pptx 二进制文件完全无法对比差异。

\vspace{4pt}
\textbf{Office 的上限，是 A Click 的起点。}
}

\notebox{【预计用时：1 分 30 秒】}

%%
%% SLIDE 7 — ADVANTAGES A-D
%%
\slidesec{7}{8 大优点（上）A--D}

\notebox{【翻到 Slide 7。扫视四个卡片。】}

\speakbox{%
我把工作流的优点归纳为八个，这里先讲前四个。

\vspace{4pt}
\textbf{A，学术严谨性。} XeLaTeX 生成的 PDF 是学术级别的排版，不是商业模板的堆砌。每次都是同样的输出，没有随机性。

\textbf{B，汇报自由度。} 叙事弧设计内置在工作流里——Hook、Roadmap、内容、综合、讨论。TikZ 图形全部代码生成，想加什么加什么。

\textbf{C，完整流水线。} 8 个阶段，没有遗漏。所有 .tex 文件可以 git 追踪，随时回滚。

\textbf{D，Plan 模式主动问需求。} AI 不猜你想要什么，先用结构化问卷问清楚，防止方向跑偏后的大量返工。
}

\notebox{【预计用时：1 分钟】}

%%
%% SLIDE 8 — ADVANTAGES E-H
%%
\slidesec{8}{8 大优点（下）E--H}

\notebox{【翻到 Slide 8。】}

\speakbox{%
后四个优点同样重要。

\vspace{4pt}
\textbf{E，三重 AI 审查。} 三个专责 agent 并行工作：slide-auditor 检查布局溢出，tikz-reviewer 检查图形碰撞，pedagogy-reviewer 检查叙事节奏。并行审查，串行修复。

\textbf{F，讲稿同步输出。} slides.pdf 和 script.pdf 永远同步交付。讲稿格式化为口播内容加舞台提示，打印出来直接能读。

\textbf{G，AI 配图自动化。} 摄入阶段就规划好配图位置，AI 生成后自动嵌入，无需用户手动找图处理分辨率。

\textbf{H，风格可配置。} 通过 ppt-intake 问卷选好字体和主色，AI 自动写入 config.tex，全局生效。换整套主题色只需改一行。
}

\notebox{【预计用时：1 分钟】}

%%
%% SLIDE 9 — PLAN MODE INTAKE
%%
\slidesec{9}{S2 Plan 模式摄入}

\notebox{【翻到 Slide 9。】}

\speakbox{%
让我详细说一下触发流程。

当你对 Cursor AI 说「帮我做个汇报」，工作流自动切换到 Plan 模式，然后用 AskQuestion 工具发来一份结构化问卷，涵盖五个维度：内容主题、格式要求、叙事风格、输出配置，以及——这是本次更新新增的——视觉风格。

视觉风格问卷问你三件事：主色系偏好（深蓝、深绿、深紫，或自定义 RGB）；中文字体（楷体、黑体、宋体）；英文字体（Times、Helvetica、Palatino）。

你回答完，AI 自动把这些写进 config.tex，一处修改，全局生效。
}

\notebox{【预计用时：45 秒】}

%%
%% SLIDE 10 — CONTENT GENERATION
%%
\slidesec{10}{自动内容生成}

\notebox{【翻到 Slide 10。】}

\speakbox{%
内容生成阶段是整个流水线的核心。

AI 从你提供的 PDF 或 DOCX 中提取关键概念和论证结构，然后逐页生成双语内容——每个 block、alertblock、exampleblock 都有对应的中英文。同时生成 TikZ 图形代码，规划配图位置。

有几个技术约定很重要：引号用 Unicode 弯引号，不用 ASCII 直引号，因为 CJK 字体环境下直引号会有字距问题；列内 TikZ 节点用 \textbackslash linewidth 不是 \textbackslash textwidth；内容溢出先减文字，必要时用 shrink 参数；每次至少编译两遍。
}

\notebox{【预计用时：45 秒】}

%%
%% SLIDE 11 — TRIPLE AI REVIEW
%%
\slidesec{11}{AI 三重质量审查}

\notebox{【翻到 Slide 11。指向 TikZ 三角图。】}

\speakbox{%
三重审查是工作流的质量门禁。

三个 agent 分工明确：slide-auditor 专门看布局——文字有没有溢出、字号是否合适、图片引用是否有效；tikz-reviewer 专门看 TikZ 图形——节点有没有碰撞重叠、标签是否可读；pedagogy-reviewer 专门看叙事——弧度是否清晰、术语是否一致、节奏是否适合听众。

三个人并行看，发现问题后串行修复，最多三轮。这比人工单轮检查的漏检率低很多。
}

\notebox{【预计用时：45 秒】}

%%
%% SLIDE 12 — TECH SPEC
%%
\slidesec{12}{S3 技术规范：字体与配色}

\notebox{【翻到 Slide 12。指向配色方块。】}

\speakbox{%
我们来看一下技术规范。

配色体系受 Pedro Sant'Anna（Vanderbilt / Microsoft Research）的 Beamer 排版风格启发：深蓝作为主色承载标题和结构，砖红作为强调色用于警告块，金色专门用于关键引言，绿色系用于示例块。

字体矩阵：中文用楷体 Kaiti SC，英文用 Times New Roman，标题加粗，术语斜体标注，代码等宽。

所有这些都在 config.tex 里集中管理，与模板本体解耦。
}

\notebox{【预计用时：45 秒】}

%%
%% SLIDE 13 — CONFIG.TEX
%%
\slidesec{13}{用户配置 config.tex}

\notebox{【翻到 Slide 13。指向左侧代码块。】}

\speakbox{%
config.tex 是整个工作流中用户唯一需要主动编辑的文件。

它做三件事：定义中英文字体名，定义主色系 RGB 值——预设了三套，你选一个取消注释就好——以及定义页脚文字。

slides-template.tex 顶部有一行 \textbackslash input\{config.tex\}，所有配置在编译时自动注入。你不需要动模板本体的任何一行。

从现在开始，你的每一套汇报都可以有自己的 config.tex，像配置文件一样管理风格，而不是每次在幻灯片里逐帧调整。
}

\notebox{【预计用时：45 秒】}

%%
%% SLIDE 14 — SCRIPT GENERATION
%%
\slidesec{14}{讲稿生成机制}

\notebox{【翻到 Slide 14。指向左侧两个示例框。】}

\speakbox{%
讲稿生成是第七阶段，在审查通过之后。

讲稿格式分两种框：蓝色框是口播内容，就是你实际要说的话，按幻灯片逐页组织；灰色框是舞台提示，包括「展示 Slide X」「停顿两秒」「指向右边图片」这类非口播的提示。

讲稿编译成 A4 PDF，页眉有标题和讲者名字，每页都有页码。打印出来，在旁边记几个时间点，就是一份可以直接用的讲稿。

你现在看到的这份讲稿，就是按这个格式生成的。
}

\notebox{【预计用时：45 秒】}

%%
%% SLIDE 15 — BEFORE vs AFTER
%%
\slidesec{15}{S4 效果对比 Before vs After}

\notebox{【翻到 Slide 15。浏览对比表格。】}

\speakbox{%
最后来看效果对比。

制作时间：传统方式 5 到 8 小时，A Click 工作流 20 到 40 分钟，节省约 85\%。

排版质量：传统方式字体乱、截图模糊、跨平台不一致；A Click 是 XeLaTeX 矢量输出，完美。

讲稿：传统方式另写，常常跟 PPT 脱节；A Click 自动生成，严格对应。

质量审查：传统方式上台前才发现问题；A Click 三个 AI agent 提前把关。

可复现性：传统方式每次手动，差异很大；A Click 的 .tex 源文件在 git 里，一键复现。
}

\notebox{【预计用时：1 分钟】}

%%
%% SLIDE 16 — DELIVERY SPEC
%%
\slidesec{16}{输出交付规范}

\notebox{【翻到 Slide 16。】}

\speakbox{%
工作流有一个硬性规定：在以下条件未满足之前，不得声明任务完成。

必须交付 slides.pdf 和 script.pdf，两者缺一不可。必须通过质量门禁：PDF 可正常渲染，无文字溢出超过 10pt，所有图片引用有效，讲稿页数等于幻灯片页数。

这个规定写在 AGENTS.md 里，Cursor 在启动工作区时会自动读取。AI agent 无法绕过它。
}

\notebox{【预计用时：30 秒】}

%%
%% SLIDE 17 — CREDITS
%%
\slidesec{17}{设计灵感与致谢}

\notebox{【翻到 Slide 17。】}

\speakbox{%
最后是致谢。

这个工作流的幻灯片排版风格参考了 Pedro Sant'Anna 在 Vanderbilt 和 Microsoft Research 时使用的 Beamer 模板规范——配色体系和幻灯片节奏特别有参考价值。内容密度和信息层级的设计参考了 Isaiah Andrews 在 Harvard 的学术汇报惯例。框架和宏设计基于 Till Tantau 等人编写的 Beamer User Guide。

这三份参考资料都是公开可得的，推荐任何想认真做学术汇报的同学去看。

链接和完整文档在 GitHub：\texttt{github.com/KyrieWong7/graduate-ppt-rescuer}
}

\notebox{【预计用时：45 秒】}

%%
%% SLIDE 18 — ENDING
%%
\slidesec{18}{结尾}

\notebox{【翻到 Slide 18（结尾页）。等大家看完屏幕上的文字。】}

\speakbox{%
「One sentence. One click. Two PDFs.」

这是这个项目的核心承诺。你提供思考，AI 负责排版。你把时间花在研究上，不是花在逐帧拖字号上。

感谢大家的聆听。这个项目完全开源，欢迎使用、fork、提 issue。

如果你有任何问题，现在或者之后都可以找我。谢谢。

\vspace{4pt}
\textit{Thank you. The project is open source at the link on screen.\\
Questions are welcome, now or later.}
}

\notebox{【鞠躬，等待提问。总计约 10 分 30 秒。】}

\vspace{20pt}
\begin{center}
  \textcolor{gold}{\rule{0.5\textwidth}{0.4pt}}\\[6pt]
  {\small\color{speakblue}\textit{本讲稿由 A Click is All You Need 工作流自动生成。\\
  This script was generated by the pipeline it describes.}}
\end{center}

\end{document}
